% \usepackage[margin=2.2cm]{geometry}
% \usepackage[T1]{fontenc}
% \usepackage{lmodern}
% \usepackage{xcolor}
% \usepackage{listings}
% \usepackage[colorlinks=true,linkcolor=blue!50!black]{hyperref}
% 
\definecolor{cmt}{HTML}{5A7A5A}
\definecolor{kw}{HTML}{1F4E9C}
\definecolor{bg}{HTML}{F6F6F6}

\lstdefinestyle{base}{
  basicstyle=\ttfamily\small,
  backgroundcolor=\color{bg},
  commentstyle=\color{cmt}\itshape,
  keywordstyle=\color{kw}\bfseries,
  frame=single, rulecolor=\color{black!20},
  breaklines=true, columns=fullflexible,
  showstringspaces=false, tabsize=4,
  xleftmargin=4pt, xrightmargin=4pt,
  aboveskip=8pt, belowskip=8pt
}
\lstdefinelanguage{gas}{
  morekeywords={pushq,popq,movq,ret,call,jmp,ud2,sti,hlt,iretq,cli},
  morekeywords=[2]{.text,.globl,.code64},
  sensitive=false, morecomment=[l]{\#}
}
\lstdefinestyle{cstyle}{style=base,language=C,
  morekeywords={uint64_t,uint32_t,size_t,_Static_assert,offsetof}}
\lstdefinestyle{asmstyle}{style=base,language=gas}
\newcommand{\code}[1]{\texttt{#1}}

%-------------------------------------------------------------
\section{Overview}
The design follows the xv6 model. The switch saves only
\emph{callee-saved} registers on the old task's kernel stack, because the
C caller already treats every other register as clobbered across the call.
The saved stack pointer lives in the task structure.

\begin{itemize}
  \item \code{struct context}: the stack image produced by the pushes.
  \item \code{struct task}: kernel stack, CR3, state, and two list nodes.
  \item \code{switch\_context}: assembly stack swap.
  \item Intrusive lists: run queue, wait queues, and all-tasks list.
\end{itemize}

%-------------------------------------------------------------
\section{Intrusive Lists}
The link node is embedded in the object, so enqueueing never allocates.
This matters in scheduler and interrupt context. A self-linked node means
``on no list''.

\begin{lstlisting}[style=cstyle]
#include <stddef.h>
#include <stdint.h>

struct list_node { struct list_node *prev, *next; };

static inline void list_init(struct list_node *h)
{ h->prev = h->next = h; }

static inline int list_empty(const struct list_node *h)
{ return h->next == h; }

static inline void list_add_tail(struct list_node *n, struct list_node *h)
{
    n->prev = h->prev; n->next = h;
    h->prev->next = n; h->prev = n;
}

static inline void list_del(struct list_node *n)
{
    n->prev->next = n->next;
    n->next->prev = n->prev;
    n->prev = n->next = n; /* self-linked = not on any list */
}

#define container_of(p, T, m) \
    ((T *)((char *)(p) - offsetof(T, m)))
\end{lstlisting}

%-------------------------------------------------------------
\section{Context and Task Structures}
\subsection{Context}
The layout must match the push order in \code{switch.S}. Lowest address
first, so this is exactly the stack image after the pushes. The \code{rip}
field is the return address consumed by \code{ret}.

\begin{lstlisting}[style=cstyle]
struct context {
    uint64_t r15, r14, r13, r12, rbx, rbp;
    uint64_t rip;
};

_Static_assert(sizeof(struct context) == 7 * 8, "context layout");
\end{lstlisting}

\subsection{Task}
\code{link} is used by the run queue \emph{or} one wait queue, never both.
\code{all\_link} stays linked for the task's whole lifetime.

\begin{lstlisting}[style=cstyle]
enum { TASK_RUNNING, TASK_BLOCKED, TASK_ZOMBIE };

struct task {
    struct context *ctx; /* saved kernel RSP; MUST be offset 0 */
    uint64_t kstack_top; /* for TSS.rsp0 */
    uint64_t cr3; /* physical address of PML4 */
    int pid;
    int state;
    struct list_node link; /* runq OR a wait queue */
    struct list_node all_link; /* all-tasks list */
};

_Static_assert(offsetof(struct task, ctx) == 0, "ctx must be first");
\end{lstlisting}

%-------------------------------------------------------------
\section{Context Switch in Assembly}
Arguments: \code{rdi} = \code{struct context **old},
\code{rsi} = \code{struct context *new}. Usage:
\code{switch\_context(\&prev->ctx, next->ctx);}

\begin{lstlisting}[style=asmstyle]
    .text
    .globl switch_context
switch_context:
    pushq %rbp
    pushq %rbx
    pushq %r12
    pushq %r13
    pushq %r14
    pushq %r15
    movq %rsp, (%rdi) # *old = current stack pointer
    movq %rsi, %rsp # switch stacks
    popq %r15
    popq %r14
    popq %r13
    popq %r12
    popq %rbx
    popq %rbp
    ret # pops ctx->rip: resumes the new task
\end{lstlisting}

\paragraph{Why it works.} \code{call} already pushed the return address,
so the six pushes complete a \code{struct context} with \code{rip} on top.
For a task that has run before, \code{ret} returns to where it called
\code{switch\_context}. For a new task, it returns into the trampoline.

%-------------------------------------------------------------
\section{Task Creation and Trampoline}
\code{switch\_context} restores only callee-saved registers, and
\code{rdi} is not one of them. So \code{entry} and \code{arg} are stashed
in \code{r12} and \code{r13} of the initial context, and the trampoline
moves them into place. \code{r12} and \code{r13} survive
\code{call task\_first\_run} because they are callee-saved.

\begin{lstlisting}[style=asmstyle]
    .globl task_trampoline
task_trampoline:
    call task_first_run # C: release sched lock, enable interrupts
    movq %r13, %rdi # arg
    call *%r12 # entry(arg)
    call task_exit # never returns
    ud2
\end{lstlisting}

\begin{lstlisting}[style=cstyle]
extern void task_trampoline(void);

void task_init(struct task *t, void *kstack_base, uint64_t size,
               void (*entry)(void *), void *arg)
{
    uint64_t top = ((uint64_t)kstack_base + size) & ~0xFULL;
    t->kstack_top = top;

    struct context *c = (struct context *)(top - sizeof(*c));
    *c = (struct context){0};
    c->r12 = (uint64_t)entry;
    c->r13 = (uint64_t)arg;
    c->rip = (uint64_t)task_trampoline;
    t->ctx = c;

    list_init(&t->link);
    list_init(&t->all_link);
}
\end{lstlisting}

\paragraph{Alignment.} The context sits directly below \code{top}, which
is 16-byte aligned. After the six pops and \code{ret}, \code{rsp == top}.
Each \code{call} in the trampoline then pushes 8 bytes, so the callee sees
\code{rsp \% 16 == 8}, as the ABI requires. A wrong value shows up later
as a crash on an aligned SSE store in compiled code.

%-------------------------------------------------------------
\section{Address-Space Switch}
Load CR3 before calling \code{switch\_context}, with interrupts disabled.
The kernel image, direct map, and all kernel stacks must be mapped at the
same virtual addresses in every address space (identical upper-half PML4
entries), or the code after the CR3 write will fault.

\begin{lstlisting}[style=cstyle]
static inline uint64_t read_cr3(void)
{
    uint64_t v;
    __asm__ volatile("movq %%cr3, %0" : "=r"(v));
    return v;
}

static inline void write_cr3(uint64_t v)
{
    __asm__ volatile("movq %0, %%cr3" :: "r"(v) : "memory");
}

extern void tss_set_rsp0(uint64_t rsp0);
extern void switch_context(struct context **old, struct context *new);

void sched_switch(struct task *prev, struct task *next)
{
    tss_set_rsp0(next->kstack_top); /* traps from user mode */
    if (next->cr3 != read_cr3()) /* avoid needless TLB flush */
        write_cr3(next->cr3);
    switch_context(&prev->ctx, next->ctx);
}
\end{lstlisting}

%-------------------------------------------------------------
\section{Run Queue and Scheduler}
Round-robin. The caller must hold the scheduler lock or have interrupts
disabled. Invariant: \code{current} is on no list, and a task is on at
most one list through \code{link}.

\begin{lstlisting}[style=cstyle]
static struct list_node runq; /* list_init(&runq) at boot */
struct task *current;
static struct task idle_task; /* never on runq */

void schedule(void)
{
    struct task *prev = current, *next;

    if (prev->state == TASK_RUNNING && prev != &idle_task)
        list_add_tail(&prev->link, &runq);

    if (list_empty(&runq)) {
        next = &idle_task;
    } else {
        next = container_of(runq.next, struct task, link);
        list_del(&next->link);
    }

    if (next == prev)
        return;
    current = next;
    sched_switch(prev, next);
}
\end{lstlisting}

\paragraph{Blocking and waking.} To block: set \code{state = TASK\_BLOCKED},
add the task to a wait queue (reusing \code{link}), call \code{schedule()}.
To wake: remove it from the wait queue, set \code{TASK\_RUNNING}, append to
\code{runq}. Enqueue and set \code{BLOCKED} under the same lock the waker
takes, otherwise a wakeup can be lost and the task sleeps forever.

%-------------------------------------------------------------
\section{All-Tasks List}
Used for PID lookup, debugging dumps, and parent/child bookkeeping.

\begin{lstlisting}[style=cstyle]
static struct list_node all_tasks; /* list_init(&all_tasks) at boot */

void task_register(struct task *t) /* caller holds task-list lock */
{
    list_add_tail(&t->all_link, &all_tasks);
}

void task_unregister(struct task *t) /* when the zombie is reaped */
{
    list_del(&t->all_link);
}

struct task *task_find(int pid) /* caller holds the lock */
{
    struct list_node *n;
    for (n = all_tasks.next; n != &all_tasks; n = n->next) {
        struct task *t = container_of(n, struct task, all_link);
        if (t->pid == pid)
            return t;
    }
    return NULL;
}
\end{lstlisting}

%-------------------------------------------------------------
\section{Invariants and Pitfalls}
\begin{itemize}
  \item Every existing task is on \code{all\_tasks} exactly once. Registering
        twice creates a cycle and scans loop forever.
  \item A task is on at most one of \{run queue, a wait queue\}.
  \item All list operations need interrupts disabled (single core) or an
        interrupt-safe spinlock (SMP).
  \item Do not free an entry while scanning. Read \code{next} before deleting.
  \item A pointer from \code{task\_find} is valid only while the lock is held.
        On SMP, add reference counting.
  \item A task cannot free its own kernel stack. Mark it \code{TASK\_ZOMBIE},
        switch away, and let a reaper free it.
  \item The idle task loops on \code{sti; hlt}.
  \item Not handled here: FPU/SIMD state (build with
        \code{-mgeneral-regs-only} or add \code{xsave}), FS/GS base MSRs,
        and per-CPU \code{current} for SMP.
\end{itemize}
\begin{lstlisting}[style=cstyle]
#include <stddef.h>

struct list_node { struct list_node *prev, *next; };

static inline void list_init(struct list_node *h) { h->prev = h->next = h; }
static inline int list_empty(const struct list_node *h) { return h->next == h; }

static inline void list_add_tail(struct list_node *n, struct list_node *h)
{
    n->prev = h->prev; n->next = h;
    h->prev->next = n; h->prev = n;
}

static inline void list_del(struct list_node *n)
{
    n->prev->next = n->next; n->next->prev = n->prev;
    n->prev = n->next = n; /* self-linked = "not on any list" */
}

#define container_of(p, T, m) ((T *)((char *)(p) - offsetof(T, m)))

enum { TASK_RUNNING, TASK_BLOCKED, TASK_ZOMBIE };

struct task {
    struct context *ctx; /* still offset 0 */
    uint64_t kstack_top;
    uint64_t cr3;
    int pid;
    int state;
    struct list_node link; /* run queue OR a wait queue, never both */
};

static struct list_node runq; /* list_init(&runq) at boot */
struct task *current;
static struct task idle_task; /* never on runq; runs when runq is empty */

/* Caller holds the scheduler lock / has interrupts disabled. */
void schedule(void)
{
    struct task *prev = current, *next;

    /* Still runnable (yield or preemption)? Put it at the back. */
    if (prev->state == TASK_RUNNING && prev != &idle_task)
        list_add_tail(&prev->link, &runq);

    if (list_empty(&runq)) {
        next = &idle_task;
    } else {
        next = container_of(runq.next, struct task, link);
        list_del(&next->link);
    }

    if (next == prev)
        return;
    current = next;
    sched_switch(prev, next); /* from the earlier message */
}
\end{lstlisting}
% \end{document}
